\begin{abstract}
We present a comprehensive theoretical framework in which information density—rather than conventional matter or energy—constitutes the fundamental substrate of physical reality. The Unified Information-Density Theory (UIDT) introduces a scalar field $S(x)$ representing information density, from which mass, gravitational interactions, spacetime curvature, and cosmological expansion emerge through a single dimensionless invariant $\gamma_{\mathrm{UIDT}} \approx 16.339$. This framework provides a constructive, non-perturbative solution to the Yang--Mills existence and mass gap problem, one of the Clay Mathematics Institute Millennium Prize Problems, predicting a mass gap $\Delta_{\mathrm{UIDT}} = 1.710 \pm 0.015~\mathrm{GeV}$ in remarkable agreement with lattice quantum chromodynamics calculations.

The theory integrates the latest observational datasets as of December 2025, including DESI Data Release~2 baryon acoustic oscillations across 14.3 million galaxies and quasars, Euclid Quick Release~1 observations of 2,674 dwarf galaxy candidates, JWST tip-of-the-red-giant-branch distance measurements, and ACT Data Release~6 cosmic microwave background polarization maps. Through the implementation of Barrow--Tsallis hybrid entropy—unifying fractal quantum gravity corrections with non-extensive thermodynamic statistics—we derive a dynamical description of dark energy evolution that resolves the persistent tension between early- and late-universe measurements.

Our key empirical validations include: (1) Dark energy equation of state at $z = 0.5$: $w(z) = -0.920 \pm 0.003$, representing a $6.4\sigma$ deviation from $\Lambda$CDM; (2) Hubble constant: $H_0 = 72.6 \pm 1.8~\mathrm{km\,s^{-1}\,Mpc^{-1}}$, reconciling early--late universe measurements to $0.1\sigma$; (3) Matter fluctuation amplitude: $S_8 = 0.756 \pm 0.002$, resolving the discrepancy between weak lensing and CMB observations; (4) Yang--Mills mass gap consistent with lattice QCD at the 99\% confidence level; (5) Predicted Casimir force anomaly of $+0.64\%$ at $d = 0.66~\mathrm{nm}$, testable in ongoing NIST experiments.

The theory yields falsifiable predictions for upcoming observational campaigns scheduled between 2026 and 2028, including measurements from the Square Kilometre Array hydrogen intensity mapping ($w(z = 1.5) = -1.03 \pm 0.0015$), final Euclid data releases, and high-luminosity LHC searches for the predicted scalar glueball resonance at $1.705 \pm 0.015~\mathrm{GeV}$. With a global reduced chi-squared of $\chi^2/\mathrm{dof} = 0.811$, representing a $>15\sigma$ improvement over $\Lambda$CDM, UIDT offers a unified, predictive, and empirically testable framework addressing fundamental questions in quantum field theory, cosmology, and the nature of dark sector physics.
\end{abstract}
\keywords{dark matter --- dark energy --- cosmological parameters --- Yang--Mills theory --- information theory --- quantum field theory --- baryon acoustic oscillations}
